\documentclass[11pt,a4paper]{article}
\usepackage[T1]{fontenc}
\usepackage[utf8]{inputenc}
\usepackage{amsmath,amssymb,xcolor}
\usepackage[margin=25mm]{geometry}
\usepackage[colorlinks=true,urlcolor=blue]{hyperref}
\definecolor{nexorablue}{HTML}{2459B5}
\setlength{\parindent}{0pt}
\setlength{\parskip}{7pt}
\setlength{\emergencystretch}{2em}
\renewcommand{\familydefault}{\sfdefault}
\begin{document}

{\small\bfseries\color{nexorablue} NEXORA CPGE}\par
{\small MPSI / NEX-CIR-04 / Énoncé}\par
{\LARGE\bfseries Énergie reçue sous deux impulsions}\par
\vspace{0.25em}\hrule\vspace{0.6em}

Un dipôle est parcouru par un courant de $30\,\mathrm{mA}$ orienté de $A$ vers $B$ pendant $0{,}20\,\mathrm s$, avec $U_{AB}=4{,}0\,\mathrm V$. Pendant les $0{,}10\,\mathrm s$ suivantes, $U_{AB}=4{,}0\,\mathrm V$ et le courant signé vaut $-10\,\mathrm{mA}$. Calculer la charge algébrique transférée de $A$ vers $B$, la puissance reçue sur chaque intervalle et l’énergie nette reçue.


\end{document}
